\documentclass[twocolumn,twocolappendix]{aastex631}
\usepackage{amsmath}




\newcommand{\vdag}{(v)^\dagger}
\newcommand\aastex{AAS\TeX}
\newcommand\latex{La\TeX}


\begin{document}

\title{ASTRO6530 Project}



\author{Jason Sevilla}





\section{VLA}
First download data.  Copy ms to new ms in same directory, but named by the transient.  Then, copy the same ms to a different directory which will be working directory.  Finally, split ms, noting the band and whether you are taking raw or calibrated data

Calibrated Data is using pipeline for CASA 6.5.4-9|2023.1.0.125

use robust = 0.5

\subsection{AT2021ahuo}

\subsection{AT2022abfc}

\subsubsection{20221223}
field 0: Secondary J0453-2807
field 1: ZTF22abvrxkk
field 2: Primary 0521+166=3C138

X-band: 0~1, 50~81 to 0~33
Ku-band spw 2~49 to 0~47



\subsubsection{20230624}
field 0: Secondary J0453-2807
field 1: AT2022abfc
field 2: Primary 3C147

flag scan 5,32

Ku-band spw 2~49 to 0~47
X-band: 0~1, 50~81 to 0~33
S-band: 82~97 to 0~15
\subsubsection{20240202}
field 0: Secondary J0453-2807
field 1: ZTF22abvrxkk
field 2: Primary 0521+166=3C138
\subsection{AT2023fhn}

\subsubsection{20230825}
field 0: J0956+2515
field 1: AT2023fhn
field 2: 1331+305=3C285

flag scan 5, 26

X-band: 0~1, 50~81 to 0~33
Ku-band 2~49 to 0 ~47
\subsection{AT2023hkw}

\subsubsection{20230614}
field 0: Secondary Calibrator J1035+5628
field 1: Target AT2023hkw
field 2: 1331+305=3C286

Use ea27 as reference, ea05, 23, 28, 12, 15 as backups
actual 0.8 for secondary, vs 0.64
ea14, ea16, spw26 problematic in B0

On 2nd file, backup for flag of short scans, and then flag of antennas, then flag spw 41

\subsubsection{20230704}
field 1: Secondary J0135+5628
field 2: AT2023hkw
field 3: 1331+305=3C286

X-band: spw 0~1
\subsubsection{20230826}
Antenna 14 excluded, Ant 12 affected by band switch problem for Ku-band, Ant 16 bad sensitivity from warm Ku-band.

C-band: spw 2~33, 0~31 after split
X-band: 0~1, 34~65, 0~33 after
Ku: 66~113, 0~47 after
field 0: Secondary Calibrator J1035+5628
field 1: Target AT2023hkw
field 2: 1331+305=3C286

Use Ref ant ea07, alternative 28

X-band: actual 0.8 Jy, get 0.6 Jy

flagged ea 21 from high gain phase on Jun 20 10:39

Ku-band : actual 0.4, get 0.41

C-band: 
Flagging: look at scan 30, spw=6, field=0

spw 17, 19 weird on bandpass gain, low on some antenna, discont corrected amplitudes
ea22 \& ea7 baseline high gain amp
ea03 has high delays

no data on spw 14, filtered or flagged out?

resolution 0.23 arcsec from max uv of 900000, want cell size of 0.05 arcsec, pick 360, image size of 432 = 21 arcsec 

im6 is best attempt, with mostly default settings

check automated pipeline,

-corr3: has a second iterative cycle, im7, doesn't help

attempting to use provided calibrated data doesn't remove it

Shami mentioned likely phase cal error

new file: \_custom in 2nd folder is raw data with custom calibration steps.  Already quack, clip, and online flagged
cal1 is first calibration attempt
used briggs weighting, much better!  first try, 5e-5 Jy, actually 7-8e-5
Second try with RFI flagging works pretty well!


\subsection{AT2023vth}
\subsubsection{20240509}
Done with no antpos

X-band: spw 0~33
C-band: spw 34~65, now 0~31
Ku-band: spw 66~113, now 0~47
Ka-band: spw 114~177, now 0~63
field 0: primary 1331+305=3C286
field 1: secondary J1751+0939
field 2: AT2023vth

ea 5, 17, 20 lost
ea 2,8,22 have newly updated baselines may be wrong
Use ea04 as reference, alternatives ea14, then 11,12,23

toss scan 8, only 8 seconds long; 24 and 25 also a few seconds long

X-band: get 2.4 Jy vs actual 2 Jy
C-band:
2 Jy vs actual 1.8 Jy for secondary, close 

ku-band:
only scan 43 for secondary calibrator
have to exclude scan 42 from flagging

ea14,25,27 a little strange on g0all phase

ea13, 27 have insane delays, flag on June 14, 3:03pm
also ea28, june 14, 7:24pm on Ku-rawth

split into file without raw

640000 wavelengths means resolution of 0.3 arcsec, want cell size of 0.05 arcsec, pick 360
get 2.9 Jy for secondary instead of 4 Jy
Looks like just noise, best is natural weighting, 2.1e-5


Ka-band: Also noise
%\begin{figure}[!h]\label{fig:beam}
% \centering
%  \includegraphics[width=\linewidth]{figures/beam.png}
%\caption{Diagram from \citet{rybicki} showing the beaming of %synchrotron radiation as the electron gyrates in a magnetic %field.  The beaming is into a half-opening angle of 1/$\gamma$.  %The beaming is visible to a faraway observer for an angle %$2/\gamma$ of the particle's circular path.}
%\end{figure}

%\begin{deluxetable*}{cccccc}
%
%\setlength{\tabcolsep}{0.25em}
%\tablewidth{0.94\textwidth}
%\tabletypesize{\footnotesize}
%\tablewidth{0pt}
%\tablecaption{Welch's t-test p-values for each set of parameters and %each $T_{eff}$ bin \label{tab:table1}}
%\tablecolumns{6}
%\tablehead{
%\colhead{$T_{eff}$ bin}
%& \colhead{FUV/NUV ratio range/mean} & \colhead{FUV standard %error/mean}
%& \colhead{NUV standard error/mean} & \colhead{Fractional FUV %Luminosity}
%& \colhead{Fractional NUV Luminosity} \\
%\colhead{(Kelvin)}
%& \colhead{FUV/NUV ratio range/mean} & \colhead{FUV standard %error/mean}
%& \colhead{NUV standard error/mean} & \colhead{Fractional FUV %Luminosity}
%& \colhead{Fractional NUV Luminosity} 
%}
%\startdata
%\vspace{0.2cm}  4800-5400 K & 0.64 & 0.98 & 4e-10 & 1e-05 & 1e-05 \\ %\vspace{0.2cm}
%4800-5400 K & 0.64 & 0.98 & 4e-10 & 1e-05 & 1e-05 \\
%4800-5400 K & 0.64 & 0.98 & 4e-10 & 1e-05 & 1e-05 \\
%\enddata
%
%\vspace{2.5mm}
       
%\tablecomments{F}
%\end{deluxetable*}

\section{Radio Collage Notes}
\subsection{AT2023fhn}
paper currently doesn't include first non-detection at t = 12, flux < 54 uJy

plot currently has t = 74 instead of t = 66 and t = 150 to t= 117

added t = 442 point

\subsection{AT2023hkw}
Changed 86 uJy to 77 in plot for t = 44, changed t=60 to t=64, changed 83 uJy to 85 in t = 117

\subsection{AT2023vth}
Changed t = 35, 43, 89 to 30, 41, 87
Addd t = 118, 204 days

\subsection{AT2022abfc}
Changed t = 430 to 438, changed 36 to 38 in plot for t=32 days.


\bibliography{mybib}
\bibliographystyle{aasjournal}

%% This command is needed to show the entire author+affiliation list when
%% the collaboration and author truncation commands are used.  It has to
%% go at the end of the manuscript.
%\allauthors

%% Include this line if you are using the \added, \replaced, \deleted
%% commands to see a summary list of all changes at the end of the article.
%\listofchanges

\end{document}

% End of file `sample631.tex'.